\documentclass[11pt,a4paper]{article}
\usepackage[utf8]{inputenc}
\usepackage[T1]{fontenc}
\usepackage{lmodern}
\usepackage[margin=2.15cm]{geometry}
\usepackage{amsmath,amssymb,mathtools,amsthm,booktabs,longtable,array,enumitem,xcolor,hyperref,listings,stmaryrd}
\hypersetup{colorlinks=true,linkcolor=blue,urlcolor=blue,citecolor=blue}
\emergencystretch=2em
\setlist{leftmargin=1.55em,itemsep=2pt,topsep=4pt}
\newcommand{\code}[1]{\path{#1}}
\newcommand{\State}[1]{\textsf{#1}}
\newcommand{\Role}[1]{\textbf{#1}}
\newcommand{\trace}{\mathcal{T}}
\newcommand{\Obj}{\mathcal{O}}
\newcommand{\Comp}{\mathcal{C}}
\newcommand{\Mem}{\mathcal{M}}
\title{Formalisme architectural de FulleGenBrain\_2\\
\large Architecture compositionnelle, autonome, vérifiable et résiliente}
\author{Spécification intégrée -- évolution de FulleGenBrain\_1}
\date{25 octobre 2026}
\begin{document}
\maketitle
\begin{abstract}
FulleGenBrain\_2 est spécifié comme un système cognitif dont le principe structurel central est la programmation compositionnelle. Les compétences, méthodes, règles, opérateurs, modèles du monde et procédures de validation sont décrits comme des objets typés, versionnés, composables et révisables. L'architecture articule observation multi-vues, empreintes DDD, maillage MG/MGR/MGR2/ASCR, synthèse de programmes et de règles, reproduction/clonage, recherche active, raisonnement spécialisé, sémantique, mémoire de provenance, contrôle autonome et vérification indépendante. Le système n'intègre et ne dépend d'aucun modèle de langage neuronal, y compris eLLM. Les techniques de représentation multi-résolution, de prédiction, de calibration ou de génération de candidats ne sont reprises que sous forme de mécanismes indépendants, explicites et évaluables. Le document distingue les contrats théoriques, les preuves conditionnelles, les sources historiques, les tests locaux et les capacités qui requièrent encore une implémentation intégrée. Il ne préjuge ni le résultat de la construction ni les performances futures.
\end{abstract}
\newpage\setcounter{tocdepth}{1}\tableofcontents\newpage

\section{Conventions et statut des affirmations}
\subsection{Version et mandat}
Ce document remplace la spécification architecturale FulleGenBrain\_1 par une version qui recentre toutes les capacités autour d'un langage de composition explicite. La théorie autonome de la programmation compositionnelle est livrée dans le même dossier et constitue le fondement théorique de la présente architecture. Les identifiants de composants, contrats et transitions ci-dessous sont des contrats cibles, pas une affirmation que toutes les parties ont déjà été implémentées.

Les marqueurs de statut sont courts et distincts :
\begin{itemize}
\item \Role{SOURCE} : constat présent dans une source examinée;
\item \Role{SPÉCIFICATION} : contrat défini par ce document;
\item \Role{PREUVE SOUS HYPOTHÈSES} : conclusion mathématique dérivée d'hypothèses explicites;
\item \Role{TEST LOCAL} : vérification sur un cas borné;
\item \Role{NON INTÉGRÉ} : chaîne complète non exécutée dans cette livraison;
\item \Role{À ÉVALUER} : question demandant une mesure sans présumer son résultat.
\end{itemize}

\subsection{Règles d'interprétation}
Un exemple qui fonctionne n'établit pas une loi universelle. Un compilateur qui réussit à construire un artefact ne prouve pas que sa sémantique métier est correcte. Une dérivation valide dépend de ses prémisses. Une bonne imitation n'est pas à elle seule une reconstruction de mécanisme. Un score heuristique n'est pas une certification. Toute sortie conserve son type, ses hypothèses, ses preuves, ses versions, son périmètre de validité et le statut qui lui est effectivement justifié.

\section{Objectif et principes non négociables}
FulleGenBrain\_2 doit pouvoir observer, représenter, organiser, composer, tester, apprendre, expliquer, reproduire, prédire, raisonner, vérifier, réviser et transformer ses propres programmes de résolution. Il ne suffit pas d'avoir les modules : il faut que leur coopération soit définie par des contrats de flux, de contrôle, de preuve, de coût et de récupération.

\begin{enumerate}
\item \textbf{Composition comme colonne vertébrale.} Toute capacité complexe est soit un composant typé, soit une composition de composants, soit un mécanisme du runtime qui vérifie et gouverne ces compositions.
\item \textbf{Aucune dépendance à un modèle de langage neuronal.} Aucun LLM, eLLM, poids de modèle de langage externe ou service de génération neuronal n'est un composant d'exécution, un fournisseur de connaissances implicite, une voie de repli ou une dépendance obligatoire.
\item \textbf{Distinction des rôles épistémiques.} L'observation, l'hypothèse, la règle candidate, la dérivation, le résultat testé et le résultat certifié sont des statuts différents.
\item \textbf{Pluralité avant fermeture prématurée.} Lorsque les données n'identifient pas un modèle unique, les alternatives et leurs désaccords sont conservés; le contrôleur peut demander une observation discriminante.
\item \textbf{Évolution tracée et réversible.} Une nouvelle compétence, une transformation ou une optimisation est un candidat soumis à typage, test, preuve disponible, évaluation et retour arrière.
\item \textbf{Résilience explicite.} Panne, timeout, manque de données, corruption, conflit et incertitude possèdent des états et stratégies propres. Aucune dégradation n'est masquée sous un statut de succès.
\item \textbf{Pas de tâche abandonnée par défaut.} Si une fonction nécessaire ne possède pas encore de mécanisme suffisant, elle génère une obligation de formalisation, de construction et d'évaluation. Le contrôleur peut exposer un état de travail incomplet sans requalifier l'objectif en impossible.
\end{enumerate}

\section{Architecture d'ensemble : un organisme à trois plans}
\subsection{Plan structurel : catalogue et compositions}
Le catalogue contient des compétences atomiques, des méthodes composées, des règles/invariants et des modèles versionnés. Chaque composant possède un contrat d'entrée/sortie, des effets, des ressources, des préconditions, des postconditions, des tests, une provenance et une politique de révocation. Les compositions sont des AST ou graphes typés qui sélectionnent les composants, raccordent leurs valeurs, organisent leurs branches, définissent les récupérations et décrivent le flux de preuves.

\subsection{Plan dynamique : runtime et contrôleur}
Le runtime exécute les compositions approuvées; le contrôleur choisit les objectifs, budgets, probes, chemins alternatifs et conditions d'arrêt; l'ordonnanceur gère la séquence et le parallélisme lorsque les effets l'autorisent. Le contrôleur ne peut pas transformer une proposition en résultat certifié. Le runtime n'invente pas de compétences absentes du registre. Le registre, le synthétiseur et le vérificateur restent des responsabilités distinctes.

\subsection{Plan évolutif : découverte et révision}
Le système propose de nouveaux blocs, règles, programmes, modèles de transition et compositions à partir d'observations, de contre-exemples, d'analogies calculables, de recherche grammaticale et de tests. Le candidat est exécuté en environnement d'essai; la validation compare la version candidate à une référence, à des données tenues à l'écart et à une suite de régression. L'activation est versionnée. Une régression ouvre une transaction de retour à la dernière version approuvée et préserve les preuves de l'épisode.

\subsection{Cycle de vie global}
\[
\begin{aligned}
\text{Objectif} &\to \text{Composer une voie} \to \text{Observer/Calculer}\\
&\to \text{Comparer les modèles} \to \text{Vérifier} \to \text{Mémoriser/Réviser}.
\end{aligned}
\]
Ce cycle n'est pas un pipeline figé. Une observation inattendue peut relancer la définition de tâche; un contre-exemple peut changer un modèle; une panne peut déclencher un composant de remplacement; une vérification invalide peut révoquer les descendants de la preuve concernée. Chaque transition enregistre sa cause, son coût, ses entrées, sa sortie et son statut.

\section{Contrat universel des objets et statuts}
\subsection{Objet commun}
Tout élément manipulable est encapsulé dans
\[
\begin{aligned}
o=(&id,kind,payload,type,context,time,source,status,scope,version,\\
   &deps,uncertainty,certificate).
\end{aligned}
\]
Le champ \code{kind} distingue donnée brute, observation, token/symbole, parse, référence, fait, règle, candidat, programme, composition, état, action, événement, trajectoire, prédiction, résultat, contre-exemple, probe, certificat et entrée mémoire. Le champ \code{source} décrit la provenance; \code{scope} précise le domaine où l'élément s'applique; \code{deps} est un ensemble de dépendances avec relations typées; \code{uncertainty} décrit les ambiguïtés restantes. Un type non encore reconnu doit être conservé comme donnée opaque versionnée, puis donné lieu à un adaptateur formalisé.

\subsection{États épistémiques}
\begin{longtable}{p{0.22\textwidth}p{0.69\textwidth}}
\toprule Statut & Signification opérationnelle \\
\midrule
\State{RAW} & donnée reçue, contenu non encore interprété \\
\State{OBSERVED} & donnée observée sous un protocole précisé avec provenance \\
\State{PARSED} & structure calculée par un parseur déclaré; pas une garantie de vérité sémantique \\
\State{HYPOTHESIS} & interprétation ou mécanisme candidat \\
\State{RULE-CANDIDATE} & règle reifiée avec portée et conditions de test \\
\State{PROGRAM-CANDIDATE} & programme ou composition de résolution proposée \\
\State{SUPPORTED} & compatible avec un ensemble explicite de preuves et un domaine précis \\
\State{DERIVED} & conséquence d'une règle ou d'un opérateur nommé et de prémisses tracées \\
\State{TESTED} & évalué sur un protocole/test identifié avec résultats enregistrés \\
\State{CERTIFIED} & accepté par un vérificateur défini dans un périmètre défini \\
\State{REFUTED} & contredit par un contre-exemple dans les conditions indiquées \\
\State{UNRESOLVED} & plusieurs candidats ou états demeurent compatibles \\
\State{STALE} & dépendance, contexte, source ou version ont changé; réévaluation obligatoire \\
\State{WITHDRAWN} & retrait explicite; interdiction d'usage comme prémisse active non revalidée \\
\State{PARTIAL} & résultat utile mais incomplet, avec les obligations non satisfaites \\
\State{FAILED} & exécution échouée avec catégorie, trace et récupération éventuelle \\
\bottomrule
\end{longtable}
Les promotions et rétrogradations sont des transactions journalisées. Une hypothèse n'est pas promue par son score seul; un candidat testé demeure lié au protocole; un certificat est limité aux hypothèses, version de vérificateur et domaine notés sur celui-ci.

\section{Programmation compositionnelle : noyau central de FulleGenBrain\_2}
\subsection{Contrat d'un composant}
Chaque bloc $c$ suit le contrat
\[
c=(id,version,A,B,Pre,Post,Effects,Resources,TracePolicy,Verifier,Recovery,Tests),
\]
où $A,B$ sont les types d'entrée/sortie, $Pre/Post$ les contrats logiques, $Effects$ le système d'effets, $Resources$ le budget, $TracePolicy$ les événements conservés, $Verifier$ le contrôle indépendant pertinent et $Recovery$ la stratégie de reprise/compensation. Les composants peuvent être purs, étatful, observationnels, transactionnels ou interactifs; ces classes sont explicites. Le contrat distingue ce qui est spécifié, testé, prouvé et mesuré.

\subsection{Grammaire conceptuelle du DSL}
La grammaire exécutable évolue par versions; la base architecturale comprend les opérateurs suivants :
\begin{align*}
p ::= {}& \mathsf{Id}\mid\mathsf{Call}(c,e)\mid\mathsf{Seq}(p_1,\dots,p_n)\mid\mathsf{Par}(p_1,\dots,p_n)\\
&\mid\mathsf{If}(g,p_t,p_f)\mid\mathsf{Map}(p,e)\mid\mathsf{Fold}(p,e,a)\mid\mathsf{Loop}(g,p,I)\\
&\mid\mathsf{Select}(p_1,\ldots,p_n)\mid\mathsf{Compose}(m,\theta)\mid\mathsf{Probe}(q)\mid\mathsf{Observe}(q)\\
&\mid\mathsf{Verify}(p,v)\mid\mathsf{Recover}(p,r)\mid\mathsf{Checkpoint}(p,k)\mid\mathsf{Txn}(p,\kappa)\\
&\mid\mathsf{Revoke}(d)\mid\mathsf{Synthesize}(G,D)\mid\mathsf{Transform}(p,\tau).
\end{align*}
La grammaire n'est pas fermée à jamais : un nouveau cas d'utilisation requiert la spécification formelle de son opérateur, l'analyse de ses effets, une sémantique, une stratégie de validation et des tests de non-régression. La compatibilité est versionnée; un runtime qui ne comprend pas un nœud ne doit pas l'ignorer silencieusement.

\subsection{Lois de base de composition}
La séquence $g\circ f$ est admissible si la sortie de $f$ satisfait le type et la précondition de $g$. L'associativité est garantie pour les compositions pures selon l'égalité des fonctions; pour les programmes avec effets, elle n'est autorisée que lorsque les traces, transactions, exceptions et effets observables restent équivalents. L'identité ne supprime aucun événement exigé par le contrat. Le parallélisme requiert une analyse de lectures/écritures, d'effets externes et de ressources. Une transformation telle que fusion, permutation, factorisation ou mise en cache est un candidat avec une obligation d'équivalence, pas une optimisation présumée sûre.

\subsection{Synthèse comme capacité de premier rang}
Le synthétiseur reçoit un objectif explicite, un ensemble d'exemples, une grammaire de candidats, des coûts, des contraintes et un oracle indépendant. Il génère plusieurs compositions, les normalise, contrôle types et effets, exécute les cas observés, recherche des contre-exemples et teste les compositions inédites. Il garde les alternatives compatibles. Le classement peut privilégier coût, fidélité, robustesse ou information gagnée, mais aucun score heuristique n'a le droit de certifier la solution.

\subsection{Apprentissage des blocs}
Lorsqu'une séquence d'actions réussit sur plusieurs exemples, le système peut proposer de la factoriser comme sous-composition nommée. Cette induction comprend : comparaison des traces; alignement des entrées/sorties; séparation des invariants et des paramètres; construction d'une abstraction; tests sur variantes; contrôle des effets; mesure de couverture; validation du contrat; ajout au catalogue comme candidat; approbation ou rejet motivé. Une macro apprise ne devient pas une primitive réputée universelle sans les tests requis.

\subsection{Réutilisation à plusieurs niveaux}
La composition s'applique à tous les niveaux : opérations élémentaires; méthodes de domaine; stratégies; workflows; contrôleurs; procédures d'apprentissage; protocoles de vérification; plans de récupération; transformations du modèle lui-même. Une compétence peut être répétée dans plusieurs voies si cela augmente la fiabilité ou donne un contrôle indépendant, à condition que les rôles et les responsabilités soient différents ou que l'indépendance recherchée soit mesurée. La duplication non justifiée des calculs est évitée par cache traçable, pas par suppression aveugle.

\section{Registre de compétences et architecture stratégique}
\subsection{Familles de blocs}
\begin{longtable}{p{0.23\textwidth}p{0.34\textwidth}p{0.32\textwidth}}
\toprule Famille & Responsabilité & Exemples de blocs \\
\midrule
Observation & capturer des entrées accessibles et établir leur provenance & observation fichier/flux/trace; parseur de données; capture d'épisode; sonde \\
Représentation & construire les vues et structures typées & tokenisation déterministe; analyse syntaxique; DDD; normalisation; types et unités \\
Maillage & relier faits, règles, dépendances, candidats et contextes & MG; MGR/MGR2; ASCR; graphe temporel; provenance \\
Synthèse & proposer et raffiner des règles, programmes et compositions & recherche grammaticale; anti-unification; synthèse d'AST; apprentissage de sous-programmes \\
Exécution & appliquer des règles, calculs, outils et compositions & eGenBrain v4 adapté; solveurs; moteur de contraintes; runtime compositionnel \\
Expérimentation & sélectionner les observations utiles et contraster des hypothèses & choix de probe; plan d'essais; mutation de cas; générateur de contre-exemples \\
Monde et trajectoires & prédire et simuler les transitions d'état & modèle de transition; planification; recherche de trajectoire; simulateur \\
Vérification & contrôler indépendamment le résultat et les preuves & vérificateur de contrat; vérifieur de preuve; contrôle d'unités; oracle de test externe \\
Mémoire & préserver données, programmes, preuves et dépendances & mémoire persistante; cache de résultats; checkpoints; registre de versions \\
Résilience & diagnostiquer et reprendre après panne & retry policy; checkpoint/restart; compensation; gestion de circuit; rollback \\
Contrôle & choisir le prochain objectif et les budgets & contrôleur de tâche; ordonnanceur; gestion de conflits; politique d'abstention \\
\bottomrule
\end{longtable}
Aucune famille ne remplace toutes les autres. La composition est l'organe d'intégration, le contrôleur l'organe de régulation, la mémoire l'organe de continuité, le vérificateur la porte de justification, et l'observation le lien avec les données réelles.

\section{Observation, empreintes DDD et sondes actives}
\subsection{Observations comme mesures}
Soit $x$ une situation et $q$ une sonde dans un contexte $c$ à une résolution $r$. La valeur observée est $O(q,x,c,r)$ et l'empreinte
\[
\phi_{Q,C,R}(x)=\big(O(q_1,x,c_1,r_1),\ldots,O(q_m,x,c_m,r_m)\big).
\]
L'empreinte est relative au jeu de sondes, au contexte et aux résolutions. Deux situations égales sous cette empreinte sont indiscernables pour ce protocole, non nécessairement identiques dans le monde. La révision de $Q,C,R$ modifie la partition observée et nécessite versionnement.

\subsection{Multi-vues et abstractions hiérarchiques}
Une observation peut être projetée dans plusieurs vues complémentaires : octets, champs structurés, tokens/symboles, syntaxe, références, sens, événements, trajectoire ou contraintes numériques. Les vues gardent une référence à la source et la transformation qui les a produites. Une vue plus abstraite doit préserver les distinctions que ses tâches aval exigent; si elle les efface, le système doit revenir à une vue plus fine ou acquérir une nouvelle observation.

\subsection{Choix de probes}
Le contrôleur maintient un ensemble de candidats $H=\{h_i\}$. Une probe est utile si ses résultats prévus séparent des candidats actuellement compatibles. La sélection peut optimiser un gain d'information estimé, le coût d'acquisition, le risque, la capacité de discrimination et la dépendance envers les modèles. Le score est un guide de recherche. La preuve de distinction découle ensuite de la sémantique de la probe et du résultat observé, lorsque les hypothèses sont suffisantes.

\subsection{Garde-fous d'observation}
Les journaux séparent observation directe, dérivation, simulation et contrefactuel. Toute observation reste dans le périmètre d’accès déclaré. Une donnée inaccessible n’est jamais enregistrée comme capturée. Si une capacité manque, le système spécifie et construit un adaptateur autorisé, ou marque explicitement le champ comme non observé.

\section{Maillage relationnel : MG, MGR, MGR2, ASCR et extensions}
\subsection{MG : mémoire persistante}
MG conserve les objets durablement réutilisables : faits observés, règles candidates et acceptées, opérateurs, programmes, compositions, certificats, contre-exemples, contextes, transitions, versions et liens de provenance. Chaque entrée indique sa portée et les conditions de révocation. Les éléments temporaires n'y sont promus qu'après validation selon une politique explicitée.

\subsection{MGR/MGR2 : maillages de travail}
MGR construit un contexte relationnel local depuis une requête, une situation ou une recherche; MGR2 affine ou restructure ce contexte après les résultats intermédiaires. Ils restent transitoires et versionnés par tâche. Une hypothèse temporaire ne modifie pas la mémoire persistante par effet de bord. La promotion vers MG est une transaction contrôlée.

\subsection{ASCR : sélection des candidats et témoins}
ASCR gère la sélection de candidats, leur provenance, leurs critères d'inclusion et les éléments qui permettent de les contredire. Le score de classement reste séparé de l'état épistémique. Le pool conserve les candidats concurrents utiles et les résultats incompatibles. Le choix d'un candidat pour la prochaine expérience ne le transforme pas en candidat vrai.

\subsection{Nouveaux maillages fonctionnels}
FulleGenBrain\_2 distingue au moins quatre vues reliées : (i) le graphe de provenance et de dépendances; (ii) le graphe d'hypothèses, objections et contre-exemples; (iii) le graphe temporel/dynamique des états, événements et transitions; (iv) le graphe de composition des programmes et compétences. Ces vues partagent des identifiants typés mais gardent des relations sémantiquement distinctes. Elles peuvent être matériellement réunies dans un store relationnel commun si les contrats de requête les distinguent.

\subsection{Fermeture, invalidation et calcul dérivé}
Chaque dérivation enregistre la liste exacte des prémisses, des règles et des opérateurs. Si un fait, une règle, un opérateur ou un contexte est retiré ou modifié, le graphe de dépendance calcule tous les descendants à réévaluer. Les cycles de dépendance sont groupés en composantes fortement connexes. Une sortie dépendant d'une prémisse retirée devient \State{STALE} puis \State{WITHDRAWN} ou est reconstruite avec une preuve indépendante. Une entrée mémoire ne peut pas demeurer \State{CERTIFIED} simplement parce qu'elle était autrefois valide.

\section{Induction de règles, programmes et modèles}
\subsection{Espace général des candidats}
Le candidat général est
\[
h=(kind,syntax,contract,scope,assumptions,evidence,predictions,cost,status,trace),
\]
où \code{kind} peut être règle, opération, programme, composition, parse, interprétation sémantique, modèle de transition ou politique. La grammaire définit les formes actuellement recherchées; elle est versionnée et peut être étendue. Le système ne doit pas limiter par principe l'induction à la découverte de règles logiques.

\subsection{Cycle de découverte}
\begin{enumerate}
\item construire des exemples observés et séparer données d'apprentissage, de validation et de test;
\item produire plusieurs candidats à partir de recherche syntaxique, induction, analogies structurées, modèles d'exécution et sous-programmes connus;
\item normaliser les candidats équivalents dans le domaine tout en préservant les alternatives non prouvées équivalentes;
\item vérifier le typage, les effets et les préconditions;
\item calculer les conséquences dans un domaine déclaré;
\item rechercher contradictions, contre-exemples et situations sous-déterminantes;
\item choisir une probe qui discrimine plusieurs candidats ou une batterie de cas adversariaux;
\item exécuter la probe, mettre à jour les preuves et conserver les résultats négatifs;
\item tester sur de nouvelles compositions et environnements;
\item promouvoir, garder comme candidat, réfuter ou marquer non résolu selon les critères.
\end{enumerate}

\subsection{Pluralité, sous-détermination et complexité}
Un corpus fini peut être compatible avec de nombreuses fonctions. Toute identification dépend d'une classe d'hypothèses, d'une grammaire, de propriétés de simplicité, d'observations discriminantes ou d'hypothèses de domaine. Les principes de parcimonie servent à ordonner la recherche, non à transformer la simplicité en preuve de vérité. Si plusieurs modèles restent compatibles, ils sont conservés ou une observation plus informative est demandée.

\subsection{Correction de modèles candidats}
La validation d'un programme candidat sépare fidélité sur les exemples, respect des contraintes, performance sur holdout, comportement sous perturbation, qualité des explications, coûts, sécurité des effets et calibration. Un candidat peut être bon sur une mesure et échouer sur une autre. Le système garde un vecteur de mesures et une décision qui cite le profil d'usage visé.

\section{Observation, reproduction et clonage comportemental}
\subsection{Représentation d'une démonstration}
Une démonstration contient la cible autorisée, le protocole de capture, les entrées, contexte, états connus, actions, sorties, horodatages, événements et incertitudes. Les données d'entraînement et de test sont séparées par épisode, entité, environnement ou graine lorsqu'une proximité permettrait la mémorisation.

\subsection{Quatre niveaux distincts}
\begin{enumerate}
\item \Role{Rejeu exact :} reproduire une trace connue. Le succès prouve seulement la reproduction du cas capturé.
\item \Role{Imitation entrée-sortie :} retrouver les sorties attendues pour les entrées couvertes.
\item \Role{Modèle de transition :} reconstruire les règles d'évolution d'un état sous une action dans un domaine observé.
\item \Role{Programme/politique compositionnelle :} découvrir une procédure réutilisable dont les invariants se maintiennent dans des variantes et combinaisons inédites.
\end{enumerate}
Ces niveaux ont des mesures et seuils indépendants. Une reproduction parfaite sur entraînement ne doit pas masquer la généralisation faible sur perturbations.

\subsection{Boucle de clonage et reconstruction}
La trace observée est convertie en données typées; plusieurs explications et algorithmes candidats sont produits; le système synthétise un clone de référence borné; il mesure les écarts; il génère des entrées nouvelles et des changements de contexte; il cherche des contre-exemples; il révise le modèle; il conserve l'original et la version candidate avec des manifestes et des rapports de comparaison. La cible cachée et les sorties de test ne sont pas injectées dans le clone. Si une interface d'observation, un format ou un domaine n'est pas encore pris en charge, une compétence correspondante est spécifiée et ajoutée, avec autorisation et protocole de contrôle.

\section{Modèle du monde, trajectoires, prédiction et causalité}
\subsection{Transition d'état}
Le monde est représenté par un modèle de transition candidat
\[
\hat{T}_{\theta}(s_{t+1}\mid s_t,a_t,c_t),
\]
ou par une relation symbolique lorsque les états, actions et règles sont explicites. Le modèle porte la source des états, les conditions, l'incertitude, l'horizon et le domaine. Une transition simulée n'est pas une observation réelle. Une séquence reconstruite n'est pas une preuve causale.

\subsection{Séparer prédiction, intervention et explication}
L’observation établit des corrélations ou régularités, pas la causalité à elle seule. Une intervention modifie une action ou une condition contrôlée. Un contrefactuel décrit un contexte alternatif simulé. Les analyses rétrospectives sont marquées comme telles. Toute conclusion causale explicite les hypothèses d’identification et les interventions disponibles.

\subsection{Planification de trajectoire par composition}
Une trajectoire candidate peut être obtenue en composant modèle de transition, opérateur de coût, contraintes, estimateur d'incertitude et vérificateur de plan. Le contrôleur explore plusieurs trajectoires sous budget, demande des observations si elles changent le choix, et recalcule après tout écart significatif. Les heuristiques de recherche ordonnent l'exploration mais ne remplacent pas les contraintes dures ni la vérification.

\section{Langage et sémantique sans modèle de langage neuronal}
\subsection{Architecture du traitement linguistique}
Le traitement linguistique s'appuie sur une chaîne de composants explicites et composables : normalisation; segmentation/tokenisation déterministe; grammaires et parseurs; analyse morphologique si nécessaire; arbre syntaxique; résolution des références; rôles sémantiques; portée de la négation/quantification/modalité; temps/aspect; contexte discursif; ambiguïtés concurrentes; conséquences testables; validation par contraintes et données. Les méthodes statistiques légères sont admises comme composants documentés, mais aucun modèle génératif de langage neuronal ou service de ce type ne fait partie de l'architecture.

\subsection{Analyse concurrente}
Pour un énoncé ambigu, le système produit plusieurs analyses structurées, chacune associée aux fragments qui l'étayent, aux hypothèses, aux contextes, aux conséquences et aux tests qui permettraient de la distinguer. La désambiguïsation peut reposer sur une grammaire, le contexte, la cohérence avec les faits, les contraintes ou une question de clarification. Le choix est révisable et garde sa provenance. L'absence d'analyse suffisante déclenche le développement de la grammaire ou du parseur manquant, pas une sortie inventée.

\subsection{Compositionnalité de la sémantique}
Le sens d'une expression est calculé à partir de la structure syntaxique, du sens des éléments et des règles de composition déclarées :
\[
\llbracket f(e_1,\ldots,e_n)\rrbracket_c = \mathsf{ComposeMeaning}_f(\llbracket e_1\rrbracket_c,\ldots,\llbracket e_n\rrbracket_c).
\]
La fonction est éventuellement relationnelle lorsqu'il y a plusieurs sens compatibles. Les relations de coréférence, quantification, portée et discours ne sont pas réduites à des scores de similarité. Les interprétations produites restent des hypothèses jusqu'à validation selon leur tâche.

\subsection{Réutilisation des idées multirésolutions}
Les représentations hiérarchiques et vues multirésolutions peuvent être réutilisées si elles sont explicites, déterministes ou correctement documentées, traçables et testées pour leur tâche : chaîne brute, tokens, syntagmes, structure, références, sens, discours. On n'introduit pas un objectif prédictif de tokens en tant que dépendance neuronale du noyau; on mesure plutôt les performances de chaque représentation et ses effets sur la recherche, l'ambiguïté, le transfert compositionnel et le coût total.

\section{Raisonnement formel et place d'eGenBrain v4}
\subsection{Responsabilité spécialisée}
eGenBrain v4 reste un moteur spécialisé de maillage relationnel, d'inférence typée et d'application de règles/opérateurs déclarés. Il traite les prémisses, hyperrègles, contraintes de types et unités, clôtures définies et traces de dérivation dans le périmètre explicitement pris en charge. Il n'est pas responsable, à lui seul, de découvrir chaque règle absente, de sélectionner toutes les observations utiles ni de certifier la vérité des prémisses.

\subsection{Intégration compositionnelle}
La composition FulleGenBrain\_2 appelle eGenBrain v4 par un adaptateur typé contenant : le contexte de requête, les prémisses admissibles, les règles/operators activés, la limite de fermeture, les contraintes d'unités, les versions, le budget et le schéma de certificat. La réponse est convertie en objet \State{DERIVED} avec ses dépendances. Le vérificateur indépendant contrôle les contraintes pertinentes et l'intégrité de la dérivation avant toute promotion à \State{CERTIFIED}.

\subsection{Application, découverte et retrait}
Si une règle explicite existe, l'engine peut l'appliquer selon sa sémantique. Si elle n'existe pas, la découverte appartient au moteur de candidats et de composition. Si une prémisse est retirée, les dépendants sont réévalués. Les caches de fermeture sont invalidés selon leur clé de prémisses, règles, contexte et version d'opérateur. Les conflits sont conservés avec leurs preuves plutôt que supprimés par dernier-écrit-gagne.

\section{Mémoire, provenance et révision continue}
\subsection{Horizons de mémoire}
La mémoire comprend : (A) épisodique -- observations et essais; (B) de travail -- contexte et MGR/MGR2 temporaires; (C) structurelle -- faits, règles, programmes et compositions approuvés; (D) procédurale -- méthodes et politiques validées; (E) de preuve -- certificats, dépendances, révocations et tests; (F) de maintenance -- incidents, checkpoints, versions et opérations à réconcilier. Ces catégories sont des vues logiques, pas nécessairement des bases physiquement séparées.

\subsection{Cycle de révision}
Chaque changement de preuve ou de contexte déclenche la recherche des dépendances affectées. L'objet concerné et tous les descendants reçoivent un statut périmé/retrait, puis sont réévalués selon les contrats. La revalidation peut construire une nouvelle preuve alternative; elle ne réactive pas silencieusement l'ancien certificat. Les historiques sont conservés comme événements passés, séparés de l'ensemble de prémisses actives.

\subsection{Provenance complète}
Tout objet issu d'une composition garde l'identifiant du programme, l'empreinte du DSL normalisé, la version du compilateur/runtime, les identifiants de blocs, les paramètres, les entrées pertinentes, les résultats de chaque étape, la configuration, le budget consommé et la preuve du vérificateur. La chaîne de hash assure la détection de modifications sur les artefacts couverts mais doit être combinée à l'authentification des acteurs si une garantie d'origine est requise.

\section{Contrôleur, objectifs et allocation de ressources}
\subsection{Objectif vectoriel}
Les mesures de tâche restent distinctes :
\[
\begin{aligned}
\mathbf{J}=(&L_{observation},L_{induction},L_{composition},L_{imitation},\\
           &L_{transition},L_{semantique},L_{contraintes},L_{calibration},\\
           &C_{calcul},R_{effets}).
\end{aligned}
\]
Ce vecteur est une structure d'évaluation, pas une fonction d'entraînement unique. Le contrôleur applique un ordre de priorités explicite : exigences dures de validité et de sécurité; préconditions; contrat demandé; objectifs qualité; coût et préférences. Un score moyen ne peut pas compenser un certificat invalide ou une entrée de provenance corrompue.

\subsection{Contrôle anytime et budget}
Chaque requête définit ou hérite d'un budget pour étapes, temps, mémoire, probes, profondeur de composition, recherche de programmes, simulations et vérification. Le budget s'étend dynamiquement si une politique autorisée le permet, ou se répartit en fonction du gain attendu et des obligations encore ouvertes. À l'échéance, le système retourne le meilleur état justifié avec statut partiel, la liste des tâches restantes, le checkpoint et le plan de reprise. L'épuisement de budget n'est pas classé comme absence de capacité.

\subsection{Sélection de voie}
Le contrôleur choisit entre exécution d'une composition certifiée, recherche de composition nouvelle, observation supplémentaire, raisonnement formel, récupération mémoire, simulation, clonage, réparation du programme ou demande d'information. Le choix utilise la disponibilité et les preuves pertinentes, l'incertitude, la valeur des informations, le coût, les risques et la réversibilité. Lorsque deux voies indépendantes peuvent vérifier un résultat, le contrôleur peut demander une double vérification; la diversité est alors documentée.

\section{Résilience opérationnelle et continuité}
\subsection{États de panne et actions}
Les pannes sont typées : entrée invalide; bloc absent; contrat violé; dépendance/version incompatible; autorisation manquante; timeout; épuisement de ressource; composant indisponible; conflit; corruption; régression; résultat indéterminé; interruption du processus. Pour chaque type, le composant déclare retry, backoff, fallback compatible, checkpoint, compensation, rollback, réconciliation, continuation partielle ou arrêt contrôlé.

\subsection{Sauvegarde, restauration et reprise}
Le checkpoint comprend état sérialisé, version de schéma, version des compositions, journal des événements confirmés, sorties intermédiaires, graphes de dépendance, opérations externes en attente, budgets restants et hashes des artefacts. Au redémarrage, les schémas et manifestes sont validés. Les opérations d'effet externe dont l'issue est inconnue passent en réconciliation plutôt qu'être exécutées une seconde fois sans contrôle. Les écritures critiques sont atomiques ou protégées par un protocole de reprise explicite.

\subsection{Tolérance au changement de composants}
Si un bloc est absent ou déprécié, le catalogue cherche un substitut dont le contrat couvre l'usage. L'équivalence est une obligation à démontrer, pas un simple appariement de noms. Si le substitut n'offre qu'une partie des garanties, le contrôleur compose une stratégie alternative qui rétablit les garanties manquantes, par exemple en ajoutant une vérification ou une observation. Si aucune voie déjà disponible n'est adéquate, le bloc manquant devient un chantier de développement, priorisé et traçé.

\subsection{Auto-réparation}
L'auto-réparation génère un candidat de correction à partir du symptôme, du journal, de la composition et du test de régression. Elle reproduit le problème dans une sandbox, formule des hypothèses, construit un patch ou une transformation, exécute les tests négatifs et positifs, compare au comportement approuvé, puis soumet l'artefact au contrôle d'intégrité. Aucun patch ne modifie en place la version active sans checkpoint et voie de retour arrière.

\section{Sécurité des effets sans limitation artificielle des tâches}
La sécurité repose sur les capacités explicites, le principe du privilège minimal, la validation des chemins, le contrôle des identifiants de composants, la non-exécution de fragments arbitraires issus d'une composition, l'isolation et la traçabilité. Une tâche qui nécessite une capacité légitime additionnelle déclenche la définition d'une nouvelle interface et une autorisation appropriée; elle n'est ni exécutée par contournement silencieux ni abandonnée sans analyse. Le langage reste évolutif via extension revue, versionnée et testée.

Les effets sont classés comme lecture, écriture, réseau, exécution, observation interactive, horloge, aléatoire, commande externe ou modification persistante. Les contrats indiquent les droits, la portée, les compensations et les conditions d'activation. Le runtime refuse une composition mal formée ou insuffisamment autorisée avec un diagnostic précis; le système peut ensuite proposer une solution alternative sûre et complète, dans les limites réelles de ses capacités.

\section{Preuve de composition et portes de vérification}
\subsection{Contrat de bout en bout}
Pour une composition $p=c_n\circ\cdots\circ c_1$, la certification exige : contrats d'entrée/sortie composables; préconditions satisfaites; effets autorisés; versions connues; trace complète; invariants préservés; résultat vérifié par le contrôleur de qualité du domaine. Sous ces hypothèses, les règles de composition de la théorie compositionnelle propagent les garanties des blocs. Le certificat nomme les hypothèses et la portée; il ne prétend pas garantir les contextes non couverts.

\subsection{Vérificateurs indépendants}
Un vérificateur de syntaxe/compilation, un vérificateur de contrat, un vérificateur métier/domaine, un contrôleur de provenance et un évaluateur de généralisation ont des rôles distincts. Plusieurs peuvent être regroupés dans une distribution logicielle mais leurs responsabilités et résultats restent séparés. Le composant qui cherche une solution ne doit pas décider seul qu'il a réussi.

\subsection{Statuts de sortie}
Les sorties admissibles sont : \State{CERTIFIED}; \State{TESTED}; \State{SUPPORTED}; \State{PARTIAL}; \State{UNRESOLVED}; \State{REFUTED}; \State{FAILED}; \State{NEEDS-OBSERVATION}; \State{NEEDS-COMPETENCE}. Les deux derniers lancent soit une nouvelle expérience, soit un travail explicite de construction et intégration de la capacité nécessaire. Aucun statut n'est remplacé par du langage rassurant qui dépasse les preuves.

\section{Protocole d'évaluation de FulleGenBrain\_2}
\subsection{Suite de tests à construire}
\begin{enumerate}
\item Composition pure : identité, séquence, associativité, normalisation et égalité fonctionnelle.
\item Composition à effets : ordre d'écriture, conflit de ressources, idempotence, rollback et compensation.
\item Compilation : validateur AST, interpréteur de référence, compilateur, exécution différée et diff des sorties.
\item Synthèse : règle cachée, programme à exemples, grammaire étendue, ambiguïté, contre-exemples et compositions inédites.
\item DDD et probes : partitions d'observations, coût d'information, stabilité des vues et changement de contexte.
\item MG/MGR/MGR2/ASCR : promotion, provenance, révocation, fermeture, conflits, cache et reprise.
\item Clonage : rejeu, imitation, transitions, perturbations et transfert hors des démonstrations.
\item Sémantique : négation, portée, quantificateurs, coréférence, références, temps, paraphrases et contexte multi-tour.
\item Résilience : crash injecté, checkpoint corrompu, composant absent, version incompatible, timeout et opérations externes ambiguës.
\item Auto-transformation : changement de composant, refactorisation, régression, évaluation OOD, approbation et rollback.
\end{enumerate}

\subsection{Métriques}
Les mesures comprennent : correction du résultat; taux de violations de contrat; validité des certificats; taux de fausses certifications; taux de réfutation correcte; couverture de recherche; diversité des candidats; gain informationnel par probe; nombre d'observations; fidélité de reproduction par niveau; généralisation sur compositions inédites; exactitude sémantique; calibration/abstention; coût CPU/mémoire/latence; consommation des probes; taux de reprise; temps de restauration; régressions induites et temps d'ingénierie nécessaire à l'ajout d'une capacité.

\subsection{Comparaisons causales et ablations}
Les essais comparent : runtime seul; catalogue sans synthèse; synthèse sans probes; DDD sans recherche active; maillage sans graphe de provenance; eGenBrain v4 seul dans son domaine; composition complète; composition sans clonage; sans modèle de transition; sans revalidation; avec contrôles de secours indépendants. Les comparaisons gardent les données, budgets, splits et oracles constants. Les ablations identifient contributions et interactions, sans préjuger du signe ou de la taille de l'effet.

\subsection{Prévention des fuites et reproductibilité}
Les jeux sont partitionnés par composant, entité, environnement, template, épisode ou famille de composition, selon la menace de fuite. Les oracles sont séparés du runtime; les cas cachés ne sont pas accessibles au synthétiseur. Les résultats publient versions, graines, budgets, erreurs, tests non lancés, dépendances, hashes et conditions d'environnement. Les résultats répétés ne sont pas additionnés sans indiquer les répétitions.

\section{Plan d'intégration recommandé}
\begin{longtable}{p{0.16\textwidth}p{0.47\textwidth}p{0.29\textwidth}}
\toprule Phase & Travail principal & Porte de sortie \\
\midrule
P0 & contrats universels, schémas, statuts, registre et provenance & schémas validés et tests de sérialisation \\
P1 & DSL/AST, interpréteur de référence, runtime séquentiel et branchements & tests sémantiques et différentiel \\
P2 & vérificateur indépendant, effets, capacités, manifests et intégrité & tests négatifs et contrats \\
P3 & MG/MGR/MGR2/ASCR, révocation et dépendances & fermeture/invalidation et reprise contrôlées \\
P4 & synthèse de règles et de compositions sur tâches structurées & holdout, contre-exemples et métriques \\
P5 & observation DDD et sondes discriminantes & valeur des probes et absence de fuite \\
P6 & modèle du monde, trajectoires et clonage & perturbations et généralisation \\
P7 & analyse sémantique composable & suite linguistique structurée \\
P8 & résilience, checkpoint, compensation et auto-réparation & panne injectée et rollback \\
P9 & intégration complète et ablations & rapports de tâches end-to-end \\
\bottomrule
\end{longtable}
Ce plan ordonne les dépendances de construction sans figer les priorités scientifiques. Une tâche critique rencontrée en avance est formalisée puis branchée à la couche dont elle dépend; les portes restent nécessaires avant de déclarer une garantie.

\section{Correspondance avec FulleGenBrain\_1, AAM O94.6 et eGenBrain v4}
\begin{longtable}{p{0.22\textwidth}p{0.35\textwidth}p{0.34\textwidth}}
\toprule Élément & Apport retenu & Modification pour FulleGenBrain\_2 \\
\midrule
FulleGenBrain\_1 & contrat universel, statuts épistémiques, DDD, maillage, clonage, probes, modèles du monde et vérification & re-centrage sur composition; suppression des dépendances de modèle de langage neuronal; capacité évolutive et résilience de bout en bout \\
AAM O94.6 & catalogue de blocs, composition déclarative, compilation, traces, manifests, build/test et démonstrateurs bornés & étendre le socle aux contrats formels, effets, programme de synthèse, transaction, révocation, reprise et transformation testée \\
eGenBrain v4 & représentation relationnelle, inférence typée et application des règles déclarées dans son périmètre & l'intégrer comme moteur spécialisé invoqué par des compositions, sans lui attribuer toute induction \\
Order & recherche de trajectoires, exploration et organisation des essais & formaliser les heuristiques comme stratégies composables et mesurables, jamais comme preuves \\
AAM (hors seule composition) & observation, réutilisation de blocs, corrections, traces et pratiques de qualité issues des artefacts disponibles & chaque capacité est réécrite en contrat et validée indépendamment de la marque/source historique \\
DDD & empreintes, partitions et distinctions dépendantes des probes & relier explicitement les changements de mesure aux hypothèses, à la sélection d'expériences et à la provenance \\
Idées d’eLLM & multi-vues, hiérarchies, évaluation des séquences & réimplémentation explicite; aucun appel, poids ou dépendance LLM/eLLM \\
\bottomrule
\end{longtable}

\section{Questions ouvertes comme tâches de construction}
Les items ci-dessous sont des chantiers concrets qui doivent être fermés par formalisation, implementation et tests; ils ne sont pas des restrictions définitives de l'architecture : preuve machine de l'équivalence interpréteur/compilateur; logique de types/effets plus expressive; sémantique de la concurrence et de l'annulation; modèles de transactions externes; apprentissage automatique de nouveaux opérateurs de programme; induction de représentations plus robustes; sémantique large du langage; observabilité d'environnements hétérogènes; évaluation à long terme de l'auto-transformation; garanties d'équité, de coût et de stabilité. Chaque item reçoit un identifiant, un propriétaire de composant, des critères de sortie, des dépendances, un banc d'essai et un journal des preuves.

\section{Conclusion et statut}
FulleGenBrain\_2 est organisé autour d'une thèse d'ingénierie : les capacités cognitives doivent être représentées, reliées et transformées comme des programmes compositionnels typés, révisables et vérifiables. Le reste de l'architecture -- observation, DDD, maillage, induction, clonage, modèles du monde, raisonnement formel, sémantique, mémoire et autonomie -- devient une collection de compétences et de contrôleurs aux contrats explicites. Cette organisation permet de raisonner sur les dépendances, le réemploi, les régressions et la continuité sans confondre les statuts de preuve.

Le présent livrable est une spécification formelle architecturale produite après examen ciblé des sources disponibles et du démonstrateur de programmation compositionnelle AAM. Il ne prétend pas certifier l'intégration de tous les composants ni une performance globale. Les preuves de ce document sont déductives sous leurs hypothèses; les résultats des sources et les tests exécutés sont cités à leur portée exacte. La prochaine étape scientifique est l'implémentation progressive du noyau compositionnel, la vérification différentielle de son sémantique, puis les expériences end-to-end avec mesures et ablations préenregistrées.

\end{document}
